\documentclass[10pt,a4paper]{amsart}
\usepackage[margin=19mm]{geometry}
\usepackage{amsmath,amssymb,mathtools}
\usepackage[scr=rsfs,cal=euler]{mathalfa}
\usepackage{xfrac}
\usepackage[unicode,hidelinks,bookmarks=false]{hyperref}
\newcommand{\ie}{i.e.\ }
\newcommand{\cf}{cf.\ }
\newcommand{\resp}{respectively\ }
\newcommand{\Subsection}{\subsection}
\providecommand{\nobreakdash}{\nobreak\hbox{-}}
\setlength{\emergencystretch}{2em}
\multlinegap0pt
\usepackage{bm}
\usepackage{bbm}
\usepackage{dsfont}
\usepackage{xspace}
\usepackage{cancel}
\usepackage{mathtools}
\usepackage{amsthm}
\makeatletter
\@ifundefined{theorem}{\newtheorem{theorem}{Theorem}}{}
\@ifundefined{prop}{\newtheorem{prop}[theorem]{Proposition}}{}
\makeatother
\title[Cluster Synchronization via Graph Laplacian Eigenvectors]{Cluster Synchronization via Graph Laplacian Eigenvectors}
\author{Tobias Timofeyev, Alice Patania}
\date{Source: arXiv:2503.18978v2 (CC BY 4.0)}
\begin{document}
\maketitle

\noindent\emph{Source and licence.} Cluster Synchronization via Graph Laplacian Eigenvectors. Version arXiv:2503.18978v2 (CC BY 4.0), distributed under the Creative Commons Attribution 4.0 International licence, as recorded in the arXiv licence metadata for this exact version and shown on the abstract page https://arxiv.org/abs/2503.18978v2. Full rights notice: licenses/arxiv-2503.18978-CC-BY-4.0.txt.\par\smallskip

\noindent\textit{Editorial scope.} Counted unit: the Prop whose source label is \texttt{prop:Best\_Approx} in the article (source0.tex lines 418--420, source label \texttt{prop:Best\_Approx}), delivered together with the article's own complete original proof of it (source0.tex lines 422--462, one \texttt{proof} environment of 41 lines). No mathematics is written, repaired or re-derived: statement and proof are the article's own text line for line. Required context retained uncounted: none. Every definition, display and label of those lines is retained unchanged. Omitted from this package and not redistributed: 1--417, 421--421, 463--859. Nothing inside a retained excerpt is omitted or altered: every retained body is delivered exactly as the article's lines are written, so no omission receipt of any kind accompanies this package. Citations: the retained excerpts carry no citation command at all, so no bibliography entry of the article is transcribed and no reference is redistributed. Reference adaptation: none was needed and none was made; every reference inside the delivered unit resolves inside it, so no unresolved reference or citation is produced. Printed slips kept literal: no printed word, symbol, hypothesis or number of the counted statement or of its proof was altered. Layout and class: the article's own class is \texttt{revtex4-1} with packages including babel, graphicx, amsmath, amsthm, amssymb, fancyhdr, lastpage, multirow, color, ifthen, enumitem, stmaryrd, bbold, subcaption; the wrapper is the lane's pinned \texttt{amsart} editorial wrapper at 10pt on A4, which restates the theorem-like environments the retained text needs and adds the article's own macro definitions that the retained text needs (none), with extra wrapper packages (bm, bbm, dsfont, xspace, cancel, mathtools). The delivered numbering is the unit's own. Assets: no retained span contains a figure environment, a graphics-inclusion command, a PDF-page inclusion, a table or a TikZ picture; no image file is copied into this package, the package declares no asset dependency and no image file is redistributed with it. Licence: version arXiv:2503.18978v2 of this article is distributed under the Creative Commons Attribution 4.0 International licence, as recorded in the arXiv licence metadata for this exact version and shown on the abstract page https://arxiv.org/abs/2503.18978v2. A full rights notice is delivered at licenses/arxiv-2503.18978-CC-BY-4.0.txt. Characters: every retained excerpt is pure ASCII, so the delivered engine, XeLaTeX with its default Latin Modern text font, reports no missing character.
\begin{prop}\label{prop:Best_Approx}
    Let $G$ be a weighted graph with graph Laplacian $L$. Let $\pi$ be a vertex partition such that the equitable error $\|\epsilon\| \leq \delta$ for eigenvector-value pair $(\lambda,v)$ of $L^\pi$. Then $Pv$ is best approximated by eigenvectors with eigenvalues close to $\lambda$. For any truncated approximation, $u$ using eigenvectors with eigenvalues within $\gamma$ of $\lambda$ and assuming $m$ of these, the bound $\|Pv-u\| \leq (\delta/\gamma) \sqrt{n-m}$ will hold. 
\end{prop}

\begin{proof}
    Consider the basis of eigenvectors of $L$. The Laplacian is symmetric and therefore diagonalizable since the graph is undirected. Therefore its eigenvalues form an orthogonal basis.  Let $\{v_1,\dotsc,v_n\}$ be a set of normal representatives, forming the spectral basis of $L$.

    With the choice of partition $\pi$, we may define the indicator matrix $P$, quotient Laplacian $L^\pi$, and consequently $E=LP-PL^\pi$. Let $(\lambda,v)$ be an eigenvalue eigenvector pair of $L^\pi$. Following from above, we see that $L^\pi v=\lambda v \implies Lw= \lambda w - \epsilon$ where $w=Pv$ and $\epsilon = Ev$.
    
    Consider the expansions $w=\alpha_1 v_1+\alpha_2 v_2+\dotsc+\alpha_n v_n$ and $\epsilon=\beta_1 v_1+\beta_2 v_2+\dotsc+\beta_n v_n$. It follows
    
    
    \begin{align*}
        Lw=\lambda w -\epsilon &\implies L(\alpha_1 v_1+\dotsc+\alpha_n v_n) = \lambda(\alpha_1 v_1+\dotsc+\alpha_n v_n)- \epsilon\\
        &\implies 
        \alpha_1 \lambda_1 v_1+\dotsc+\alpha_n \lambda_n v_n = (\lambda \alpha_1 - \beta_1 ) v_1+\dotsc+ (\lambda \alpha_n -\beta_n) v_n\\
        \implies 
        0 &= (\lambda \alpha_1 - \beta_1 - \alpha_1\lambda_1 ) v_1+\dotsc+ (\lambda \alpha_n -\beta_n -\alpha_n\lambda_n) v_n\\
        &= ((\lambda-\lambda_1) \alpha_1 - \beta_1) v_1+\dotsc+ ((\lambda-\lambda_n) \alpha_n -\beta_n) v_n\\
        \implies 
        0 &= ((\lambda-\lambda_i)\alpha_i-\beta_i) \\&\hspace{20pt}\text{for all $i$, by orthogonality}
    \end{align*}
    By the assumption $\|\epsilon\|\leq \delta$ it follows $|\beta_i|\leq\delta$. 
    Suppose that $H$ is the basis transformation matrix for the eigenbasis of $L$ (its rows are the $v_i$). Then $\|\beta\|^2=\langle 
    \beta,\beta \rangle=\langle 
    H\epsilon,H \epsilon \rangle=\langle 
    \epsilon,H^TH\epsilon \rangle=\langle 
    \epsilon,\epsilon \rangle=\|\epsilon\|^2$ since $\langle \cdot,\cdot\rangle$ is a norm and $H$ is orthonormal. Thus $\|\epsilon \|=\|\beta\|$ and therefore any bound on the norm of $\epsilon$ carries over to the norm of $\beta$. Since $\|\cdot\|$ denotes the euclidean norm, it follows $|\beta_i|\leq \|\beta\|$ and therefore that $|\beta_i|\leq \delta$ for $1\leq i\leq n$
    
    Recall that $L$ is positive semidefinite, and therefore has only non negative eigenvalues. From the above result we have that $\beta_i=(\lambda-\lambda_i)\alpha_i$ for $1\leq i\leq n$. Thus $|(\lambda-\lambda_i)\alpha_i| \leq \delta$ for all $i$.
    
    It follows that when $\delta$ is small, $(\lambda-\lambda_i)$ or $\alpha_i$ is small. For a fixed $\beta_i$ value, the further away $\lambda_i$ is from $\lambda$, the smaller $\alpha_i$ must be. This means that eigenvectors of $L$ with eigenvalues close to $\lambda$ can have larger coefficients $\alpha_i$ and therefore contain the most information of $w$. Let $A=\{a_1,\cdots,a_m\}$ be the set of coefficients for eigenvalues within $\gamma$ of $\lambda$. For all other eigenvalues, $\alpha_i$ is likely small. Therefore $u= \alpha_{a_1}v_{a_1}+\cdots+\alpha_{a_m}v_{a_m}$ is our approximation of $Pv$ in the spectral basis. This allows for the derivation of a formal bound.
    \begin{align*}
        \|Pv - u\|^2 &= \|\sum_{i\not\in A} \alpha_i v_i\|^2\\
        &= \sum_{i\not\in A} |\alpha_i|^2    &\text{as shown earlier for basis change}\\
        &\leq \sum_{i\not\in A} \left(\frac{\delta}{|\lambda-\lambda_i|}\right)^2\\
        &= \delta^2 \sum_{i\not\in A} \left(\lambda-\lambda_i\right)^{-2}\\
        &\leq \delta^2 \sum_{i\not\in A} \gamma^{-2}\\
        &= (\delta/\gamma)^{2} (n-m) 
    \end{align*}
    Thus $\|Pv-u\| \leq (\delta/\gamma) \sqrt{n-m}$.

    
    This tells us that in a partition on a graph that deviates from an AEP mildly, the eigenvectors with eigenvalues close to $\lambda$ describe $Pv$ well.
\end{proof}

\end{document}
